\section{Definitions}%
\label{sec:definitions}
\begin{definition}
	Given a solution mapping $\mu$.
	A \emph{mapping tuple}, $\xi$, is an extension on the solution mapping
	$\mu$ with a special variable, $v_{fragment}$, bound to a Literal value,
	with the condition that $v_{fragment} \notin V$.
	The special variable $v_{fragment}$ is used to store the information
	about the \emph{fragment} a solution mapping belongs to.
	A \emph{fragment}, in the context of this work, is an identifier used for
	grouping the mapping tuples.
    Since mapping tuple is an extension on top of the existing definition of
    solution mappings, in Section~\ref{sec:preliminaries}, the operations
    defined in Section~\ref{sec:preliminaries} also applies to the mapping
    tuple.
	\todoi{TODO: Check the following sentence again to make sure it can be called as a \emph{relation}}
	Finally, a relation denoted as $\Xi$, is defined as a relation containing
	a multi-set of mapping tuples.
	\label{def:mapping_tuple}
\end{definition}


\newpage

\todoi{TODO: Rework the following definitions of iterators according to the
discussion}
\begin{definition}
	An \textbf{iterator}, $I$, consists of one or more \textbf{fields},
	$\phi^{I}$, an iterator \emph{path}, and a \emph{reference formulation}.
	An iterator extracts a list of records from the data source according to the
	given iterator \emph{path}~\cite{delva2021RMLFields}, while the
	\emph{reference formulation} determines the data format~\cite{dimou_ldow_2014}
	of the data source.
	This is the entry point to further extract nested data structures with \emph{fields}.
	\label{def:iterator}
\end{definition}




